%% ---------------------------------------------------------------------------
%% Context-Dependent Selection Among Established Urban Routing Policies
%% Journal manuscript.  Standalone: compile main.tex directly.
%%   pdflatex main -> bibtex main -> pdflatex main -> pdflatex main
%% Figures are external PDFs in figures/.  Tables are inline.
%% ---------------------------------------------------------------------------
\documentclass[preprint,12pt]{elsarticle}

\usepackage[T1]{fontenc}
\usepackage[utf8]{inputenc}
\usepackage{lmodern}
\usepackage{amsmath,amssymb}
\usepackage{graphicx}
\usepackage{booktabs}
\usepackage{tabularx}
\usepackage{array}
\usepackage{microtype}
\usepackage[hidelinks]{hyperref}
\usepackage{url}

\journal{Transportation Research Part C}

%% Float placement: keep tables and figures near the text that cites them.
\setcounter{topnumber}{3}
\setcounter{bottomnumber}{2}
\setcounter{totalnumber}{4}
\renewcommand{\topfraction}{0.92}
\renewcommand{\bottomfraction}{0.55}
\renewcommand{\textfraction}{0.06}
\renewcommand{\floatpagefraction}{0.68}

\newcommand{\CO}{CO\textsubscript{2}}
\newcolumntype{R}{>{\raggedright\arraybackslash}X}
\newcolumntype{d}{>{\raggedleft\arraybackslash}p{0.085\linewidth}}
\setlength{\tabcolsep}{4.5pt}

\begin{document}

\begin{frontmatter}

\title{Context-Dependent Selection Among Established Urban Routing
Policies: Complementarity, Decision Value and Generalisation Under
Distribution Shift}

\author{Anonymised for review}
\address{Affiliations withheld for review}

\begin{abstract}
A route guidance system must commit to a routing objective before any trip is
assigned, and the merit of that objective depends on the operating context.
We study the selection of the routing \emph{policy} itself as a per-instance
decision problem. Four established policies --- shortest distance, reactive
travel time, capacity-aware load balancing and reliability-aware mean--variance
guidance --- are evaluated on a 648-context full factorial over corridor demand,
green ratio, guidance penetration, information lag, disruption and alternative
capacity, with five paired seeds per cell: 12{,}960 valid microsimulation runs.
Three of the four policies are strictly best somewhere with a seed-resolved
margin, and 73.0\,\% of contexts have more than one non-dominated policy across
journey time, \CO\ and stopped delay, so the portfolio is complementary. Its
decision value is small: the best fixed policy leaves 0.45\,s per vehicle of
headroom on a 474.4\,s mean journey time, whereas the worst fixed policy costs
220.2\,s (46.4\,\%). A gradient-boosted selector trained on advantage rather than
on the winner's identity recovers 55.8\,\% of that headroom under
leave-one-context-out evaluation and halves the worst case, while a transparent
mechanistic rule is \emph{more} accurate and \emph{more} costly than doing
nothing. Across eight distribution shifts reported separately by type the
selector helps on two, is indistinguishable on three and harms on three; a
support-and-confidence gate never certified a deviation, because the conformal
interval is two orders of magnitude wider than the effect, and flagged none of
the domain shift its features cannot encode. The contribution is a measurement
rather than an algorithm: complementarity, decision value and deployability are
distinct and quantitatively disconnected here, so reporting the first as
evidence for the third overstates the case for adaptive routing-policy
selection. Results are scoped to one synthetic network.
\end{abstract}

\begin{keyword}
routing-policy selection \sep algorithm selection \sep decision regret \sep
distribution shift \sep selective prediction \sep traffic microsimulation
\end{keyword}

\end{frontmatter}

\section{Introduction}
\label{sec:intro}

Every deployed route guidance system embeds an answer to a question it is rarely
asked to justify: by what objective should its paths be computed? Minimising
distance, tracking the most recently observed travel times, balancing predicted
flow against link capacity and penalising the variance of travel time are all
defensible answers, and they are not the same answer --- each sends a given
request down a different corridor. The objective is therefore a decision
variable in its own right, logically prior to and constraining the route
computation that follows it. It is that prior decision, rather than the route
computation, that this paper measures.

The question is operational rather than algorithmic. An operator does not
usually get to invent a routing mechanism; it gets to choose among established
ones, and to change that choice as conditions change. Whether it should is not
obvious. The relative merit of routing objectives is known to depend on the
traffic state: reacting to observed travel times helps most when there is
congestion to react to, and can hurt when there is not, because a policy that
diverts traffic from the shortest corridor pays the diversion cost whether or
not the corridor is loaded. Modern systems also expose far more of the operating
context than demand alone. Guidance reaches a varying and often unknown share of
vehicles; the information it acts on is delayed by a measurable lag; signal
timing determines how much of the available capacity each movement receives;
capacity itself changes with disruption. Each of these is a plausible reason for
the preferred policy to change, and each is measurable before any vehicle
departs.

Treating an operating condition as a problem instance and a routing objective as
a solver places the question inside the per-instance algorithm-selection
framework \citep{rice1976,bischl2016aslib}, which supplies both the baselines
and the vocabulary used here. The reason for adopting that framing is not
convenience but the separation it forces. Whether some policies beat others in
some conditions, how much money that variation is actually worth, and whether
any rule can collect it from information available in advance are three
questions with three different answers, measured in three different units --- a
count of contexts, a cost per vehicle, and a held-out regret under shift. They
are routinely reported as though the first implied the last. Work on routing
objectives typically establishes the first and leaves the other two to
inference, which is safe only when the gap between them is small. Measuring that
gap is the purpose of this study.

We therefore ask:

\begin{description}
\item[RQ1] Under which combinations of demand, signal timing, guidance
penetration, information lag, disruption and alternative capacity do established
routing policies exhibit complementary strengths?
\item[RQ2] Which physical properties of the context explain where the preferred
policy changes?
\item[RQ3] How much decision value is available from context-dependent
selection, relative to the strongest fixed policy and to the hindsight bound?
\item[RQ4] Does a learned selector contribute anything beyond a transparent
mechanistic rule, and by what measure?
\item[RQ5] How does selector performance change under interpolation,
extrapolation, concept shift and domain shift, and what can an abstention
mechanism actually prevent?
\end{description}

The study answers these on a controlled microsimulation benchmark: a
three-corridor urban network, four established policies, a 648-cell full
factorial over the six factors above, five paired seeds per cell, and 12{,}960
valid runs with no failures and no vehicle teleports. The contributions are
empirical and methodological, and one of them is negative:

\begin{enumerate}
\item A controlled characterisation of routing-policy complementarity across six
operating dimensions simultaneously, with an explicit account of ties and of
which wins survive replication noise (Section~\ref{sec:complementarity}).
\item A quantitative separation of complementarity from decision value. The
portfolio is complementary on three of four policies and on the criterion
vector in 73.0\,\% of contexts, yet the headroom over the best fixed policy is
0.095\,\% of mean journey time, while the penalty for the wrong fixed policy is
46.4\,\% (Section~\ref{sec:value}).
\item Evidence that decision-oriented evaluation reverses the ordering that
accuracy would give. A depth-2 mechanistic rule is more accurate than doing
nothing and costs more; a gradient-boosted model trained on advantage is less
accurate than a logistic classifier and costs less
(Section~\ref{sec:selection}).
\item A distribution-shift analysis reported separately by shift type, in which
the selector helps on two splits, is indistinguishable on three and harms on
three, and in which the support gate flags every context of three extrapolation
splits but none of the domain shift it was intended for
(Section~\ref{sec:shift}).
\item A negative result about selective prediction in this regime: the
confidence gate never certified a deviation in any context, in or out of
distribution, because the conformal half-width exceeds the available advantage
by two orders of magnitude (Section~\ref{sec:selective}).
\end{enumerate}

Nothing here is proposed as a new method. Each policy is taken from the
literature unchanged, criteria are compared by dominance rather than through a
new preference model, and the selectors are standard classifiers and regressors
used off the shelf. All results belong to one synthetic environment and none has
been checked against field data.
Section~\ref{sec:lit} positions the study, Section~\ref{sec:problem} states the
decision problem, Sections~\ref{sec:method} and~\ref{sec:design} describe the
method and the computational study, Section~\ref{sec:results} reports results,
Section~\ref{sec:discussion} interprets them, and
Section~\ref{sec:conclusion} concludes.
